\begin{afterword}
\summaryhead

The post argues that almost every generalization about the world has exceptions, and that generalizations about what to do inherit them. Its thesis is that instrumental values often have no description that is both compact and local: both short and stated only in terms of the actions themselves. The example is a box opened by a machine controlled by a dozen keys. A short rule must describe the machine, and a rule about the keys alone must list every sequence. A rule that almost always works may still miss a key that burns the money. The post adds that the need for closure, which higher stakes can increase, makes such exceptions hard to tolerate.

It then warns against the reverse inference. Complicated actions do not show complicated goals. Killing anyone, even Hitler, is a large negative terminal utility, and shooting Hitler is still right because of the lives saved. To conclude that his death has positive utility is the type error of ``Terminal Values and Instrumental Values'', a kind of double counting. Whether utilities themselves are leaky, the post ends, would be a separate argument.

\responsehead

The post makes two points, and the second is the stronger. That complicated expected utilities do not show complicated utilities follows from the model of ``Terminal Values and Instrumental Values'', and the Hitler case applies it clearly. The post also marks the limit of the point: utilities ``might be'' leaky too, and that ``would be a separate argument.''

The first point, that instrumental values often have no compact and local description, is illustrated rather than shown. The box is built so that no simple rule about keys exists, which shows what the thesis means. That real instrumental values are ``often'' like this depends on the analogy between the machine and the world.

The claim about psychology gives one direction of an effect that runs both ways. ``Need for closure'' is Arie Kruglanski's term, and in Kruglanski's theory stakes can act in either direction. A cost of staying undecided raises the need for closure; a cost of being wrong can raise a need to avoid closure, and in the experiments Webster and Kruglanski report, fear of being evaluated reduced the effects that time pressure increased. The post's claim that higher stakes ``can'' increase the need holds for stakes of the first kind.

The opponents are unnamed. The relativist's clause about circumcision can be read as a claim that cultures hold different values, which is what the author, with ``I do not think'', said that dispute concerns in the earlier post; read so, the reply about Hitler does not address it. The ``many'' who infer from consequences that death is not always bad are not quoted. The best-known philosophers who deny that a death always counts the same way, the moral particularists, argue from a general claim about reasons, a view the post's closing concession leaves open. The author's own view, that the disvalue of a death stays fixed whatever follows, is close to W. D. Ross's ``prima facie duties,'' duties that keep their weight even when others outweigh them. In the terms of the earlier post, it treats each death as a separate, fixed term added into the utility of the whole outcome. That model does not require utilities to add up in this way, and the post's ``They might be!'' leaves the question open.

In short: a sound point that complicated actions do not prove complicated goals, together with a thesis about instrumental values supported only by a constructed example, a claim about stakes that Kruglanski's theory supports only in part, and opponents who are not named.
\end{afterword}
